% !TeX root = ../papers/nc-p3.tex
\section{二つの長さ十の点を持つ乗法群窓の二成分}
\label{sec:simple-weight-finite-windows}
\begin{theorem}\label{thm:simple-weight-finite-windows}
有限点グラフに作用する乗法群の矢重みが $0,1,2,3$ である全窓は、添字を保ち印を外すと次の二族になる。
\begin{enumerate}
\item $(u,v)\in(k\setminus\{0,1\})^2$,
$4u^2-4u+4v^2-4v+1\ne0$。異なるパラメータは全て非同型である。
一添字の巡回は向きを一つ選ぶと $(u,v)\mapsto(v,1-u)$ であり、最小周期は
$(1/2,1/2)$ で1、それ以外で4である。
\item $b\in k\setminus\{0,1\}$, $3b^2-3b+1\ne0$、
同型は $b'=b$ または $1-b$ と同値で、全て最小周期1である。
\end{enumerate}
両成分はそれぞれ二次元と一次元で、第一の $(1/2,1/2)$ と第二の $1/2$ だけが一致する。
二つの判別式の零点でも下記の表示は正則で、そこで点対応が正次元になる。
全表示は二次PBWで多重安定である。
\end{theorem}

\begin{proof}[必須係数とPBW]
基底を重み順に $0,1,2,3$ と呼び、四つの数字でテンソル係数を指定する。
前定理の点対は両極の二対だけである。重み1,2,4,5の切断は階数1で、両側の全係数が非零となる。
実際、各和にあるどの座標点対も辺に入らないので、その座標で少なくとも一関係は非零である。
重み0,6では関係がない。重み3の二関係を両端の $0,3$ の直線の包に制限すると、
$03,30$ 上への射影は同型である。階数が1以下なら二つの射影直線の積に高々一本の式しかなく、
グラフが有限でない。従って重み3切断の両側で $03,30$ を選んだ小行列は可逆で、
隣り合う $00,33$ は消えるから $0303,3030$ の非零対角行列である。
片側だけ $12,21$ を選ぶ八係数は、他方向の階数1ブロックに入り全て非零である。

以上から関係の先頭語を、左辞書順で数字の降順二文字六つにできる。重み1,2,4,5で順に
$10,20,31,32$ が得られる。重み3では $30$ を先に掃き出した後の $03$ の関係に
$21$ が非零係数で現れるため $30,21$ の二つになる。これは上の対角小行列と非零の交差係数から従う。
長さ3の標準候補は20個であり、四テンソルの四つの独立な一文字縮約から実次元も20以上である。
従って全ての三次の重なりが解消し、定理\ref{thm:split-inverse-regularity}と同じDiamond補題と
最上テンソルの完全対合で全次数AS正則となる。この議論には点グラフの有限性の代わりに
今得た非零係数と切断階数だけで十分である。
\end{proof}

\begin{proof}[六変数の再現正規形]
和6の語のうち隣り合う和が全て1以上5以下であるものを用いる。
$1212,2121$ の二係数を $x,y$ とし、残る38語の係数は全て非零である。
二切断の重み1,2,4,5の階数1は40本の二項式で与えられ、その指数行列は階数20である。
各語の四位置と数字への接続行列の階数は12である。辞書順にこの行列の独立行を最初から選び、
その12語の係数を1にする。行格子のSmith非零対角成分は全て1で、残余有限被覆はない。
40式にこの12正規化を加えて指数を行簡約すると、次の辞書順の六係数が独立変数となる。
\[
 (w_{3030},w_{3102},w_{3111},w_{3120},w_{3201},w_{3210})=(a,b,c,d,h,f).
\]
残りはこれらのLaurent単項式である。この規則と \path{computations/verify_simple_weight_torus.py} の
\texttt{monomial\_chart} は全係数を復元する。

重み3の両切断の順を $03,30,12,21$ とすると行列は順に
\[
 \begin{pmatrix}1&0&1&d/c\\0&a&1&c\\h/c&1&x&1\\b&dh/(bc)&1&y\end{pmatrix},
 \qquad
 \begin{pmatrix}a&0&1&1\\0&1&dh^2/(bcf)&b/c\\d&1&y&1\\f&1&h/b&x\end{pmatrix}.
\]
左上を掃き出して右下を零とする。最初の対角式は $x=h/c+1/a$, $y=bd/c+dh/(ab)$ を与える。
残りから
\[
 d=a(c-b)/c,\qquad h=bc^2(1-b)/(c-b),\qquad f=1+a(h-b)/c,
 \qquad (ab-c)(a(1-b)-c)=0 .
\]
従って $b,c$ は非零、$b\ne1$, $b\ne c$ である。$a=c/b$ のとき全残余式は恒等的に零。
$a=c/(1-b)$ では $2b=1$ は前枝で、他は $c=1$ が必要十分である。
実際一式の分子は単項式を除き $(2b-1)^2(c-1)(bc+b-c)$ で、最後の因子を零にすると
残る式は $b(2b-1)\ne0$ の下で非零である。
第一族の座標を $u=b$, $v=1-b/c$ とし、第二族は $c=1$ の $b$ を用いる。
この範囲で $a,b,c,d,h,f$ は全て単元であり、$x,y$ の零は許される。
\end{proof}

\begin{proof}[有限グラフと同型]
辺の二点の最低座標が一方だけ非零なら重み1,2,3の式を順に使って他方の全座標が消える。
最高座標も同じである。両方で最低・最高ともに零の場合は重み2,4で $11,22$ が消え、
重み3の交差係数から点対はない。従って最低どうしか最高どうかのアフィンチャートで尽きる。
最低座標を共に1とし第一点の重み1,2の座標を $t,z$ とする。重み1,2の式で第二点の重み1,2を解き、
重み3の二式で両方の最高座標を解く。残る重み4,5の式は
\[
 A z^2+Bzt^2+Ct^4,\qquad t(Dz^2+Ezt^2+Ft^4)
\]
で、$A\ne0$ は重み2,4の必須係数から従う。原点以外では $t\ne0$ で $t=1$ にできる。
全辺と反転した最高チャートで二式の終結式を取ると、単元を除き第一族で
$4u^2-4u+4v^2-4v+1$、第二族で $3b^2-3b+1$ となる。
終結式が零なら非固定軌道の曲線があり、非零なら四つのグラフ全てが有限である。
PBWの証明はこの因子を使わない。

有限グラフ上では前定理の階数1と極大トーラスの共役により、印を外す追加操作は重みの反転だけである。
反転後の12係数を再び1にすると第一族は恒等、第二族は $b\mapsto1-b$ となる。
同じ再正規化で一巡回は第一族で $(u,v)\mapsto(v,1-u)$、第二族で $b\mapsto1-b$ となり、
同型・交叉・周期の結論を得る。上のコードは終結式およびこの単項式座標変換も検算する。

最後に全ての語は一位置だけ違う別の語と共存しない。24の順列は全て現れ、
それらの重みのアフィン包は行和・列和を固定した全単体行列の空間である（置換行列の凸包）。
追加語の数字頻度を一様頻度から引いた差は、符号を揃えて
$A=(1,-1,-1,1),-A,C=(-1,2,-1,0),D=(0,-1,2,-1)$ で尽き、全四種が現れる。
$2A-A+C+D=0$ なので各追加語に十分小さい正の重みを与えて列和を揃えられる。
残る各位置の偏りは、24順列の一様平均が相対内部にあるためその正の係数を少し調整して打ち消せる。
よって対角群について多重安定であり、非対角モーメントが恒等零なので全群でも多重安定となる。
\end{proof}

\begin{proposition}\label{prop:double-middle-finite-pbw}
有限点グラフのもう一方の重み $0,1,1,2$ でも、前定理の関係重複度と端点の評価、
二つの隣接切断、四つの一文字全階数を持つ有限候補は全て二次PBWかつ標準Koszul左右AS正則である。
\end{proposition}
\begin{proof}
最低・最高文字を $x,z$、中央を二平面 $U_i$ とする。重み1の二関係は $xU,Ux$ の間の可逆なグラフ、
重み3の二関係も $zU,Uz$ の間の可逆なグラフである（両極の一意な点評価）。
重み2の $xz,zx$ への射影は同型である。そうでなければ両極文字の二直線の積に曲線が残る。
よって $xz=f_i, zx=g_i$, $f_i,g_i\in U_i U_{i+1}$ と書ける。
両者は二次行列として可逆である。実際、四極文字成分は $xzxz,zxzx$ の非零対角のみであり、
$f_i,g_i$ は相補位置を $xz,zx$ にした二中央文字成分の非零倍である。
その文字並びをもう一方向で切ると、重み1対3の階数2行列の二次正方部分である。
両側で二つの射影が同型なため、各部分が可逆になる。

各 $U_i^*$ に $f_i$ の零点列を一つ取る。可逆性から任意の一点より両方向へ延ばせる。
位置ごとに $u_i,v_i$ のうち大きい方 $v_i$ の双対がその点となる基底を取れば、
$f_i$ に $v_i v_{i+1}$ は現れず $v_i u_{i+1}$ は現れる。
$x<u<v<z$ の左辞書順で、重み1では $ux,vx$、重み3では $zu,zv$、
重み2では $zx,vu$ が先頭語となる。先と同じ三文字計算で全ての重なりが解消する。
\end{proof}
